\section{Estudio de factibilidad}

El proyecto cuenta con los recursos de hardware necesarios para llevar a cabo simulaciones de alto nivel computacional. El equipo de cómputo disponible permite implementar técnicas de paralelización, asegurando la capacidad de realizar múltiples inferencias en tiempos razonables. Adicionalmente, se dispone de herramientas tecnológicas de alto rendimiento que brindan la ventaja computacional requerida para ejecutar la simulación de manera eficiente.

En cuanto a la factibilidad económica, la mayor parte del software utilizado es de libre acceso, incluyendo Python, C++, Visual Studio Code y TeXstudio. Se utiliza también Wolfram Mathematica, que requiere una licencia anual de 195 dólares, la cual se encuentra disponible de manera gratuita para estudiantes del IPN, garantizando así la viabilidad financiera del proyecto.

Desde el punto de vista operativo, el diseño de la interfaz facilita la definición de reglas y la ejecución de simulaciones, optimizando el flujo de trabajo del grupo de investigadores. Esto permite acelerar significativamente el proceso de experimentación y análisis de resultados, maximizando la eficiencia del equipo durante el desarrollo del proyecto.

Finalmente, se cumple la factibilidad temporal, ya que el cronograma de desarrollo estimado permite realizar todas las simulaciones necesarias para el análisis posterior, asegurando que el proyecto se pueda completar dentro del tiempo previsto.
